\documentclass[12pt,a4paper]{article}
\usepackage{ctex}
\usepackage{graphicx}
\usepackage{geometry}
\usepackage{array}
\usepackage{setspace}
\usepackage{xeCJK}
\usepackage{indentfirst}
\usepackage{amsmath}
\usepackage{amsthm}
\usepackage{float}
\usepackage{subfigure}
\usepackage{booktabs}
\usepackage{multirow}
\usepackage{caption}
\usepackage{listings}
\usepackage{xcolor}
\usepackage{matlab-prettifier}

% 页边距
\geometry{left=2.5cm,right=2.5cm,top=2.5cm,bottom=2.5cm}
% 默认宋体小四
\setCJKmainfont[BoldFont=SimHei]{SimSun}
\zihao{-4}
\linespread{1.25} % 小四字通常 1.25 倍行距更接近 Word 默认

\lstdefinestyle{Matlab-boxed}{
  style=Matlab-editor,     % 继承你现在的配色与语法高亮
  frame=single,            % 画单线边框
  framerule=0.6pt,         % 边框粗细
  rulecolor=\color{black}, % 边框颜色
  framesep=6pt,            % 代码与边框的内边距
  backgroundcolor=\color{black!2}, % 淡灰底，可按需删掉
%   xleftmargin=2em, xrightmargin=2em % 让框整体内缩一点（可选）
}

\begin{document}
% 实验名称（小二黑体、居中）
\begin{center}
  {\zihao{2}\heiti 实验名称}
\end{center}
% 年级 专业 姓名 座位号（宋体小四、居中、段前后 0.5 行）
\vspace{0.5\baselineskip}
\begin{center}
  （\textbf{年级}\ \underline{2024 级}\quad \textbf{专业}\ \underline{x} \quad \textbf{姓名}\ \underline{Pinco Li} \quad \textbf{座位号} \ \underline{5}）
\end{center}
\vspace{0.5\baselineskip}

\section{实验目的}

% 在此填写实验目的

\section{实验仪器}

\begin{enumerate}
    \item 
    \item 
\end{enumerate}

% 如需插入仪器图片
% \begin{figure}[H]
%     \centering
%     \begin{minipage}[t]{0.48\textwidth}
%         \centering
%         \includegraphics[width=\textwidth]{assets/仪器1.jpg}
%         \caption{仪器名称1}
%         \label{fig:instrument1}
%     \end{minipage}
%     \hfill
%     \begin{minipage}[t]{0.48\textwidth}
%         \centering
%         \includegraphics[width=\textwidth]{assets/仪器2.jpg}
%         \caption{仪器名称2}
%         \label{fig:instrument2}
%     \end{minipage}
% \end{figure}

\section{实验原理}

% 在此填写实验原理

% 公式示例：
% \begin{align}
% E = mc^2 \tag{1}
% \end{align}

\subsection{子标题（如需要）}

% 子章节内容

\section{实验步骤}

% 在此填写实验步骤

\section{数据记录与处理}

% 实验数据和处理说明

\subsection{实验数据}

% 数据表格示例：
% \begin{table}[H]
% \centering
% \caption{实验数据表}
% \begin{tabular}{|c|c|c|}
% \hline
% 测量项 & 数值 & 单位 \\
% \hline
%  &  &  \\
% \hline
% \end{tabular}
% \end{table}

\subsection{数据处理}

% 数据处理过程

\subsection{实验结果图像}

% 插入图片示例：
% \begin{figure}[H]
%     \centering
%     \includegraphics[width=0.8\textwidth]{assets/结果图.png}
%     \caption{实验结果图}
%     \label{fig:result}
% \end{figure}

% 如需添加附录和MATLAB代码
% \appendix
% 
% \section{数据处理MATLAB代码}
% 
% \begin{lstlisting}[style=Matlab-boxed,caption={数据处理代码}]
% clear; clc; close all;
% % 在此编写MATLAB代码
% \end{lstlisting}

\end{document}
